\chapter{Templates}\label{app:templates}

This appendix provides three templates for direct use. They convert the formal objects of the
main text into engineering documents.

\section{Architecture Map Template}

The template records the level ladder of a system. Each row is one level.

{\small\begin{longtable}{L{1.1cm}L{2.2cm}L{2.9cm}L{2.5cm}L{2.6cm}L{2.3cm}}
\caption{Architecture map}\label{tab:tpl-architecture}\\
\toprule
\textbf{Level} & \textbf{Participant} & \textbf{Guaranteed invariants} &
\textbf{Assumed of the level below} & \textbf{May be ignored} &
\textbf{Verification method} \\
\midrule
\endfirsthead
\toprule
\textbf{Level} & \textbf{Participant} & \textbf{Guaranteed invariants} &
\textbf{Assumed of the level below} & \textbf{May be ignored} &
\textbf{Verification method} \\
\midrule
\endhead
\bottomrule
\endlastfoot
$\Lk{4}$ & & & & & \\
$\Lk{3}$ & & & & & \\
$\Lk{2}$ & & & & & \\
$\Lk{1}$ & & & & & \\
$\Lk{0}$ & & & & & \\
\end{longtable}}

The following requirements apply when filling in the table:

\begin{enumerate}[label=(\arabic*)]
\item The ``may be ignored'' column must not be left blank. A blank means the level is not a
      level (Principle \ref{pr:level-real}).
\item ``Guaranteed invariants'' must be written as propositions that can be checked, or at
      least decided.
\item ``Verification method'' must state which class of assurance in
      Table \ref{tab:assurance-matrix} the invariant falls under (mechanical, monitoring,
      consensus).
\end{enumerate}

\section{Module Contract Template}

\begin{description}
\item[Identity] Module name, the level it belongs to, unique identifier.
\item[Responsibility] What problem it solves (one sentence, no longer than one line).
\item[Boundary] What it is responsible for, and what it is explicitly not responsible for.
\item[External interface] The functions, types or interfaces through which the level above
      uses it.
\item[Precondition $P$] What the caller must satisfy.
\item[Postcondition $Q$] What is guaranteed on return.
\item[Invariant $I$] Where it is enforced; why circumvention is impossible.
\item[Side effects $E^{+}$ / $E^{-}$] What is permitted to occur; what is explicitly
      forbidden.
\item[Resource budget $R$] Upper bounds on time, space, bandwidth and concurrency, and their
      cost model.
\item[Failure semantics $F$] What happens on error, cancellation, retry and recovery.
\item[Environmental assumptions $A$] The hardware, operating system, load and external-service
      conditions on which the conclusions depend.
\item[Ownership and lifecycle] How key objects are created, initialised, closed and destroyed;
      who owns them.
\item[Concurrency model] Thread-safety guarantees, lock ordering, reentrancy.
\item[Dependencies] Which modules it depends on; whether the dependency is on an abstract
      interface or on a concrete implementation.
\item[\textbf{Oblivious set}] What this module is permitted never to care about (required; see
      Definition \ref{def:oblivious}).
\item[Verification requirements] Which tests or checks can demonstrate the above conclusions;
      whether there is an isolable unit of verification.
\item[Leak-pair risk] Known scenarios in which ``the interface is the same but the external
      behaviour differs''.
\item[Open questions] Design questions not yet settled.
\end{description}

\section{Decision Record Template}

\begin{description}
\item[Identifier and title] \code{ADR-XXXX}: decision title.
\item[Status] Draft / accepted / superseded / deprecated.
\item[Context] What situation forces the decision to be made.
\item[Decision drivers] The constraints and goals that bear on the decision (performance,
      consistency, team size, rate of evolution, cost).
\item[Candidate options] Listed one by one, each with its assumptions and its cost.
\item[Decision] Which one is chosen, and the reason in one sentence.
\item[Impact] What the decision makes easier and what it makes harder.
\item[Full argument for the excluded options] Why the other options were rejected. This item
      is the most valuable part of a decision record: it records the information and judgement
      available at the time, so that future readers need not derive it again.
\item[Verification] How the consequences of the decision are checked; which external evidence
      would show that it is wrong.
\item[Reassessment conditions] What circumstances should prompt this decision to be revisited.
\end{description}

\begin{principle}
If an architecture is a physical fact like a low-level ABI, there is no need to write a
decision record.
\emph{The ubiquity of decision records is itself evidence that high-level architecture has no
single correct answer.}
\end{principle}
